Para determinar si dos programas son contextualmente equivalentes, no basta con comparar sus valores finales --en caso de existir-- sino que también deben considerarse sus efectos sobre el estado del programa y cualquier resultado o consecuencia en la ejecución que permita distinguir un programa de otro.

Lo anterior requiere describir el comportamiento del programa en términos de cómo sus instrucciones modifican el estado del sistema durante la ejecución, aspecto abordado por la semántica dinámica del lenguaje \cite{ldp_n06}.

Entre los principales enfoques para describir el comportamiento dinámico  se encuentran la semántica operacional, la semántica denotativa y la semántica axiomática.

\subsection{Semántica operacional}
Este enfoque semántico describe paso a paso cómo se ejecuta un programa, considerando los estados intermedios además del resultado final \cite{nielson1992semantics}.

La semántica operacional se formaliza mediante un sistema de transiciones que define dos tipos de configuraciones:
\[
\begin{aligned}
\langle S, s \rangle &\quad \parbox[t]{0.8\columnwidth}{%
indica que la instrucción $S$ será ejecutada desde el estado $s$, y} \\
s &\quad \text{representa un estado final}
\end{aligned}
\]


Existen dos enfoques para especificar la semántica operacional, los cuales difieren en la manera de definir la relación de transición \cite{nielson1992semantics}. 

La relación de transición se define mediante reglas de inferencia, cuya aplicación sucesiva permite construir árboles de derivación que representan formalmente las distintas etapas de ejecución de un programa. Esto permite razonar formalmente sobre los programas y sus propiedades, ya que dichos árboles justifican los pasos de ejecución necesarios para demostrar una determinada propiedad \cite{nielson1992semantics}.

\subsubsection{Semántica natural}

También conocida como \emph{semántica de paso grande}, es un modelo de la semántica operacional que describe cómo se obtiene directamente el resultado final de la evaluación de las expresiones. 

La \enquote{semántica natural} de un lenguaje $T$ está dada por \cite{hutton2023semantics}:
\begin{itemize}
    \item un dominio $V$ de valores semánticos, y
    \item una relación de transición entre términos de $T$ y valores en $V$, que asocia a cada término los valores que pueden obtenerse mediante su ejecución completa. 
\end{itemize}

Esta relación describe cómo la ejecución de una instrucción transforma un estado inicial en un estado final \cite{nielson1992semantics}. 

Una transición en semántica de paso grande se escribe
\[
\langle S, s \rangle \Rightarrow s'
\]
para representar que la evaluación de $S$ en el estado $s$ da lugar al estado $s'$. 

Si el estado \(s'\) existe, se dice que la ejecución de \(S\) \emph{termina}; en caso contrario, se dice que la ejecución  \emph{diverge}. 


Definida la relación de transición para la semántica de paso grande, se puede definir formalmente cuándo dos programas se consideran equivalentes. 

\begin{definition}[Equivalencia en semántica natural]
Dos programas $S_1$ y $S_2$ son semánticamente equivalentes, si para todos los estados $s$ y $s'$  
\[
\langle S_1, s \rangle \Rightarrow s'
\quad \text{si y solo si} \quad
\langle S_2, s \rangle \Rightarrow s'
\]
\end{definition}


La definición anterior establece que dos programas son equivalentes si, partiendo del mismo estado inicial, terminan exactamente en los mismos estados finales.


Los estados finales dependen de la categoría sintáctica considerada. En el caso de las expresiones, la equivalencia se define por la obtención del mismo valor de evaluación. En el caso de los comandos, la equivalencia se caracteriza por producir el mismo estado final cuando se ejecutan desde un mismo estado inicial.

Además, esta definición considera equivalentes a dos programas cuando, para un mismo estado inicial,  ninguno de ellos puede derivar un valor o estado final. En consecuencia, la definición no distingue entre programas que se atascan y programas que divergen, ya que en ambos casos no existe una derivación hacia un resultado final.


%----------------------------

\subsubsection{Semántica estructural}
También conocida como \emph{semántica de paso pequeño}, su propósito es describir cómo se llevan a cabo los pasos individuales de las ejecuciones \cite{nielson1992semantics}.


La \enquote{semántica estructural} de un lenguaje $T$ está dada por \cite{hutton2023semantics}:
\begin{itemize}
    \item un dominio $S$ de estados de ejecución, y
    \item una relación de transición en $S$ que relaciona cada estado con todos los estados a los que se puede llegar realizando un solo paso de ejecución.
\end{itemize}

La relación de transición tiene la forma
\[
\langle S, s \rangle \to \gamma
\]
y describe un paso individual en la ejecución de \(S\). 

Si \(\gamma\) es de la forma:

\[
\begin{aligned}
\langle S', s' \rangle &\quad \parbox[t]{0.8\columnwidth}{%
la ejecución no ha terminado y la configuración resultante representa el cálculo restante.} \\
s' &\quad \parbox[t]{0.8\columnwidth}{la ejecución de \(S\) desde el estado \(s\) ha finalizado en el estado \(s'\).}
\end{aligned}
\]


Una configuración \(\langle S, s \rangle\) se considera \emph{atascada} si no existe  una configuración \(\gamma\) tal que $\langle S, s \rangle \to \gamma$.

Dado un programa $S$ y un estado $s$, una secuencia de derivación desde la configuración inicial $\langle S,s\rangle$ puede ser:

\begin{itemize}
    \item \emph{finita} 
    \[\gamma_0,\gamma_1,\ldots,\gamma_k\]
    tal que $\gamma_0=\langle S,s\rangle$ y
    $\gamma_i \to \gamma_{i+1}$ para todo
    $0\leq i<k$, donde $\gamma_k$ es una
    configuración terminal o atascada.

    \item \emph{infinita} 
    \[
    \gamma_0,\gamma_1,\gamma_2,\ldots
    \]
    tal que $\gamma_0=\langle S,s\rangle$ y
    $\gamma_i \to \gamma_{i+1}$ para todo
    $i\geq 0$.
\end{itemize}

La ejecución de un programa $S$ desde un estado $s$ se considera que \emph{termina} si existe una secuencia de derivación finita
desde $\langle S,s\rangle$ hasta una configuración terminal. En tal caso se denota $\langle S,s\rangle \to^{*} \gamma$.

Por otro lado, se considera que la ejecución \emph{diverge} si existe una secuencia de derivación infinita
desde $\langle S,s\rangle$.


Una vez definidas las nociones asociadas a la relación de transición, se presenta la definición formal de equivalencia en semántica de paso pequeño.

\begin{definition}[Equivalencia en semántica estructural]
Dos programas $S_1$ y $S_2$ son semánticamente equivalentes, si para todos los estados $s$

\begin{itemize}
    \item $\langle S_1, s \rangle \to ^* \gamma $ si y solo si $\langle S_2, s \rangle \to ^* \gamma$ donde $\gamma$ es una configuración atascada o terminal, y 
    \item existe una derivación infinita iniciada en $\langle S_1, s \rangle$ si y solo si existe una derivación infinita iniciada en $\langle S_2, s \rangle$. 
\end{itemize}
\end{definition}

A grandes rasgos, la equivalencia en semántica estructural es análoga a la equivalencia en semántica natural. La principal diferencia radica en que la semántica de paso pequeño describe explícitamente cada transición de la ejecución, proporcionando una caracterización más detallada del comportamiento de los programas.

Dos programas son equivalentes si, para todo estado inicial, alcanzan exactamente las mismas configuraciones terminales o atascadas mediante un número finito de transiciones. En particular, cuando ambos programas se atascan, la equivalencia se mantiene si ninguno de ellos puede continuar su ejecución mediante la aplicación de alguna regla de transición.

Asimismo, la definición contempla programas que no terminan. En este caso, dos programas son equivalentes cuando, para todo estado inicial, ninguno de ellos alcanza un valor o estado final debido a que sus ejecuciones pueden prolongarse indefinidamente.


\subsection{Semántica denotativa}
Este tipo de semántica se interesa por describir qué efecto produce la ejecución de un programa \cite{nielson1992semantics}, abstrayendo los detalles de su evaluación e implementación mediante su representación en términos matemáticos. En lugar de explicar cómo se ejecuta un programa, se enfoca en caracterizar qué significa.

La \enquote{semántica denotativa} de  un lenguaje $T$ está dada por \cite{hutton2023semantics}:
\begin{itemize}
    \item un dominio $V$ de valores semánticos, y
    \item una función semántica 
    \[
        \llbracket \cdot \rrbracket : T \to V,
    \]
    que asigna a cada término $t \in T$ un valor semántico $\llbracket t \rrbracket \in V$.
\end{itemize}
La notación $\llbracket \cdot \rrbracket$, conocida como \emph{corchetes semánticos}, denota el resultado de interpretar un término dentro de un modelo matemático del lenguaje \cite{hutton2023semantics}.

La idea es asignar a cada categoría sintáctica una función semántica que asocie cada construcción del lenguaje con un objeto matemático ---frecuentemente una función--- que describa su comportamiento. Este enfoque se basa en el principio de composicionalidad: el significado de una construcción compuesta se define únicamente en función del significado de sus subcomponentes \cite{nielson1992semantics}.

La composicionalidad garantiza que términos denotacionalmente iguales pueden sustituirse dentro de cualquier construcción sin alterar su significado, lo que justifica el uso del razonamiento ecuacional \cite{hutton2023semantics}.  Asimismo, dado que la semántica denotacional se define por inducción estructural sobre la sintaxis, en la práctica es natural emplear inducción estructural para demostrar propiedades.  %pp 3-4

A continuación, se presenta la equivalencia entre programas mediante una modelación matemática.

\begin{definition}[Equivalencia en semántica denotativa]
Dos programas $S_1$ y $S_2$ son semánticamente  equivalentes si y solo si 
\[\llbracket S_1 \rrbracket  = \llbracket S_2 \rrbracket\]
\end{definition}

Esta definición de equivalencia reduce el problema a la igualdad matemática de los valores denotados. Es decir, dos programas son equivalentes cuando sus respectivas funciones semánticas lo son, lo que implica demostrar la \textit{extensionalidad} de las mismas. 

Por lo tanto, aunque los programas posean estructuras sintácticas distintas o sigan estrategias de ejecución diferentes, ambos son indistinguibles desde el punto de vista denotacional si representan exactamente el mismo comportamiento abstracto.



\subsection{Semántica axiomática}
El enfoque axiomático define el significado de los programas en términos de especificaciones lógicas sobre su comportamiento, centrándose en las propiedades que pueden demostrarse sobre los estados del programa antes y después de su ejecución \cite{cornell2014lecture09}. La idea central es expresar afirmaciones lógicas sobre dicha ejecución. 


Una \enquote{afirmación} es una \textit{terna de Hoare} expresada de la forma
\[\{P\}\ S\ \{Q\}\]

donde 
\[
\begin{aligned}
S &\text{ denota un programa y}, \\
P, Q &\text{ son predicados lógicos.}
\end{aligned}
\]

Esta terna de \emph{correctitud parcial} expresa que toda ejecución terminante de \(S\) que parte de un estado donde se cumple la precondición \(P\), concluye en un estado que satisface la postcondición \(Q\).  

La \enquote{semántica axiomática} de un lenguaje $T$ está dada por \cite{aldrich2018hoarelogic}: 

\begin{itemize}
    \item un conjunto de axiomas sobre $T$
    \item un conjunto de reglas de inferencia que permiten establecer la veracidad de las afirmaciones. 
\end{itemize}


Este tipo de semántica puede analizarse desde dos perspectivas fundamentales de la lógica matemática: la \enquote{teoría de modelos} y la \enquote{teoría de la demostración}~\cite{Pierce:SF2}.

\subsubsection{En la teoría de modelos}

La teoría de modelos estudia las relaciones entre los lenguajes formales y las estructuras matemáticas, siendo la \enquote{verdad} el puente que conecta ambos ámbitos \cite{Manzano1989}. Desde esta perspectiva, la semántica axiomática se fundamenta en un concepto formal de \emph{validez} para las afirmaciones sobre programas.

Se dice que una terna de Hoare es \enquote{semánticamente válida} si y solo si, para todo estado $s$, se cumple que~\cite{nielson1992semantics}
\[
(s \vDash P \land \langle S,s\rangle \Rightarrow s') \implies s' \vDash Q.
\]
En tal caso, se escribe
\[
\vDash \{P\}\,S\,\{Q\}.\]

Es decir, una tripleta es \emph{válida} cuando, para cualquier estado inicial \(s\) que satisfaga la precondición \(P\), toda ejecución terminante del programa \(S\) conduce a un estado \(s'\) que satisface la postcondición \(Q\) \cite{winskel1993formal}.

Esta noción de validez se apoya en un modelo semántico basado en transiciones de estados, que describe formalmente el comportamiento de los programas \cite{Pierce:SF2}. En consecuencia, el significado de una terna de Hoare depende de cómo las ejecuciones del programa, descritas mediante transiciones de estado, satisfacen la especificación dada.


A partir de esta caracterización semántica puede definirse la equivalencia entre programas.

\begin{definition}[Equivalencia en semántica axiomática y teoría de modelos]
Dos programas $S_1$ y $S_2$ son semánticamente equivalentes si, para toda precondición $P$ y toda postcondición $Q$,
\[
\vDash \{P\} S_1 \{Q\}
\quad \text{si y solo si} \quad
\vDash \{P\} S_2 \{Q\}.
\]
\end{definition}

Desde la teoría de modelos, dos programas comparten el mismo significado si son indistinguibles ante cualquier especificación lógica; es decir, cualquier descripción válida sobre el primer programa debe ser mutuamente válida para el segundo bajo el modelo de estados subyacente. 


\subsubsection{En la teoría de la demostración}

La teoría de la demostración, o teoría de la prueba, estudia las demostraciones en los sistemas lógicos, analizando su estructura, los procedimientos para transformarlas y las traducciones sintácticas entre teorías formales \cite{TroelstraSchwichtenberg2000}.

En este sentido, la semántica axiomática se centra en la derivabilidad de las afirmaciones sobre programas. En particular, una terna de Hoare se considera válida cuando puede obtenerse a partir de los axiomas y reglas de inferencia del sistema lógico. Esto se expresa mediante la notación
\[
\vdash \{P\}\ S\ \{Q\}
\]
que indica que la terna de Hoare puede demostrarse sintácticamente.

La equivalencia axiomática de programas, en la teroría de la prueba, se define de la siguiente manera.

\begin{definition}[Equivalencia en semántica axiomática y teoría de la demostración]
Dos programas $S_1$ y $S_2$ son derivacionalmente equivalentes si, para toda precondición $P$ y toda postcondición $Q$,
\[
\vdash \{P\}\ S_1\ \{Q\}
\quad \text{si y solo si} \quad
\vdash \{P\}\ S_2\ \{Q\}.
\]
\end{definition}

A diferencia del enfoque basado en teoría de modelos, donde el significado de una terna se define mediante las ejecuciones del programa sobre estados, la teoría de la demostración interpreta las afirmaciones en términos de su derivación sintáctica. Es decir, el significado de una especificación está determinado por las reglas que permiten construir una demostración de ella \cite{Pierce:SF2}. Por ello, dos programas son equivalentes cuando satisfacen las mismas especificaciones que pueden derivarse a partir de las reglas de inferencia. 